% Assembled Team 3 project report in the supplied ELCVIA format.
% Requires elcvia.sty and the figure files listed in report/README.md.
\documentclass[twoside,a4paper,11pt]{article}
\usepackage{times}
\usepackage{elcvia}
\usepackage[english]{babel}
\usepackage{graphicx}
\usepackage{booktabs}
\usepackage{amsmath}
\usepackage[hidelinks]{hyperref}
\usepackage[capitalize,nameinlink,noabbrev]{cleveref}

\newcommand{\etal}{et al.}

\title{Trustworthy Pneumonia Detection from Chest X-rays with Grad-CAM Validation}
\author{Abhishek Singh, Ajit Yadav, Aman Singh and Himanshu Kumar\\
  \small Team 3, Section A, Project P03}
\date{}

\begin{document}
\maketitle
\pagestyle{myheadings}
\markboth{\small\itshape Abhishek Singh et al. / Trustworthy Pneumonia Detection}
         {\small\itshape Abhishek Singh et al. / Trustworthy Pneumonia Detection}
\hrulefill

\begin{abstract}
Deep neural networks can classify chest X-rays accurately while relying on non-diagnostic shortcuts such as burned-in text, laterality markers or image borders. We study this risk on the paediatric Chest X-Ray Images (Pneumonia) dataset using a custom convolutional neural network and an ImageNet-pretrained DenseNet-121. On the 624-image test set, Grad-CAM assigns a mean of 43.4\% and 39.5\% of heatmap energy to the lungs, respectively, against a mean lung area of 37.7\%; selected examples show attention on regions outside the lungs. A DenseNet-121 trained on lung-masked images achieves similar matched-test accuracy (0.8702 versus 0.8734) and higher mean lung attention (51.79\% versus 39.52\%). Cross-testing the unmasked model on masked images lowers its accuracy to 0.6314. This sensitivity to image context does not establish a causal shortcut; segmentation errors and the dataset's limits remain important caveats.

\vspace{0.3cm}
\noindent
{\em Key Words}: Pneumonia detection, chest X-rays, convolutional neural networks, DenseNet-121, Grad-CAM, shortcut learning, lung masks.
\end{abstract}
\hrulefill

% =====================================================================
% Introduction and problem statement  (author: Abhishek Singh, Stage 1)
% Prepared for the ELCVIA project report.
% =====================================================================
\section{Introduction}
\label{sec:intro}

Pneumonia is the single largest infectious cause of death in children worldwide; in 2019 it accounted for 14\% of all deaths of children under five~\cite{who_pneumonia}.
Chest X-rays are the most common imaging test for diagnosing it, but reading them takes trained radiologists,
who are scarce in many of the places where childhood pneumonia is most common.
Deep convolutional neural networks (CNNs) are therefore an attractive way to support diagnosis:
CheXNet~\cite{rajpurkar2017chexnet}, a 121-layer DenseNet, exceeded the average F1 score of four practising radiologists at detecting pneumonia
on frontal chest X-rays, and Kermany \etal~\cite{kermany2018identifying} showed that transfer learning can classify paediatric chest X-rays as
normal or pneumonia with high accuracy.

High accuracy alone does not show that a model is trustworthy. A CNN can reach a good score by learning \emph{shortcuts}~\cite{geirhos2020shortcut}:
features that happen to be correlated with the label in the training data but have nothing to do with the disease.
On chest X-rays, such shortcuts include text markers, laterality letters, image borders and hospital-specific acquisition details.
Zech \etal~\cite{zech2018variable} found that pneumonia models picked up hospital-specific information and generalised worse to other hospitals,
and DeGrave \etal~\cite{degrave2021ai} showed that COVID-19 detectors relied on markers and regions outside the lungs instead of lung findings.
A model that relies on such cues may fail without warning when it is used on images from a different scanner or hospital.

\paragraph{Problem statement.}
In this project we train a pneumonia classifier on paediatric chest X-rays and test \emph{where it looks}.
We ask two questions:
(i)~what fraction of a model's Grad-CAM~\cite{selvaraju2017gradcam} attention falls inside the lungs, and
(ii)~how removing everything outside the lungs (lung-masked training) changes accuracy and attention.

\paragraph{Approach and contributions.}
The work is organised in four stages that build on each other:
\begin{enumerate}
  \item A reproducible data pipeline with a \emph{patient-level} validation split, and a custom CNN baseline trained from scratch (\cref{sec:data}).
  \item DenseNet-121 transfer learning in the style of CheXNet, compared with the baseline.
  \item Grad-CAM heatmaps, lung masks from a pretrained segmentation model, and a lung-attention metric: the share of the heatmap inside the lung mask.
  \item Training on lung-masked images, a cross-test of the unmasked model on masked images, and a final comparison of accuracy and attention.
\end{enumerate}
  \paragraph{Main findings.}
  On the 624-image test set, the lung-masked DenseNet-121 achieved similar accuracy to the unmasked model (0.870 versus 0.873)
  and higher mean lung attention (51.79\% versus 39.52\%). When the unmasked model was evaluated on masked images, accuracy fell to 0.631.
  This indicates sensitivity to removing extra-pulmonary context, but does not identify a specific shortcut or establish causation.
Beyond the main experiments, our analysis of the raw data shows that image aspect ratio and colour mode differ between the two classes (\cref{sec:data-eda}).
Both are potential shortcuts that are present before any model is trained.


\section{Related Work}
% =====================================================================
% Related work: Abhishek's paragraph (papers used in Stage 1)
% Paste into the shared Related Work section; the others add one paragraph each.
% =====================================================================
\paragraph{Pneumonia detection from chest X-rays.}
CheXNet~\cite{rajpurkar2017chexnet} fine-tuned an ImageNet-pretrained 121-layer DenseNet on the ChestX-ray14 dataset
of more than 100,000 frontal chest X-rays and reported pneumonia detection performance that exceeded the average of four practising radiologists.
It also visualised its predictions with class activation maps, an early step towards checking where such models look.
The paediatric dataset we use was introduced by Kermany \etal~\cite{kermany2018identifying}, who applied transfer learning from ImageNet to separate normal and pneumonia X-rays,
and it is widely used as a benchmark through its Kaggle release~\cite{mooney2018kaggle}.
Later work has questioned what such models actually learn: Zech \etal~\cite{zech2018variable} found that pneumonia classifiers exploit hospital-specific signals and lose accuracy on data from other hospitals,
and DeGrave \etal~\cite{degrave2021ai} found that COVID-19 classifiers rely on markers and regions outside the lungs, examples of shortcut learning in the sense of Geirhos \etal~\cite{geirhos2020shortcut}.
Our Stage~1 builds the baseline and the data pipeline for testing exactly this on the Kermany data. Unlike much previous work on this dataset, it uses a patient-level validation split and keeps the official test set untouched.

% =====================================================================
% Related work: Ajit's paragraph (papers used in Stage 2)
% Paste into the shared Related Work section; the others add one paragraph each.
% =====================================================================
\paragraph{Transfer learning with DenseNet.}
DenseNets connect every layer to all later layers in a block, which reuses features and needs fewer parameters than comparable networks~\cite{huang2017densenet}.
Pretraining such a network on ImageNet~\cite{deng2009imagenet} and then fine-tuning it on a smaller dataset is a standard way to train with limited medical data,
because the early layers learn general features that transfer to new tasks~\cite{yosinski2014transferable}.
CheXNet~\cite{rajpurkar2017chexnet} used exactly this recipe with a 121-layer DenseNet for chest X-rays.
We follow it in Stage~2: an ImageNet-pretrained DenseNet-121 with a single output logit, first trained with a frozen backbone and then fine-tuned with a low learning rate.
Unlike CheXNet, we train on the small paediatric dataset and choose all settings on a patient-level validation set, never on the test set.

% =====================================================================
% Related work: Aman's paragraph (papers used in Stage 3)
% Paste into the shared Related Work section; the others add one paragraph each.
% =====================================================================
\paragraph{Explaining classifiers and finding shortcuts.}
Grad-CAM~\cite{selvaraju2017gradcam} highlights the image regions that raise a class score by weighting the last convolutional feature maps with their average gradients.
It needs no change to the network, which is why it is widely used to check medical image classifiers.
To turn a heatmap into a number, Score-CAM~\cite{wang2020scorecam} measures the share of the heatmap's energy that falls inside the annotated object, an \emph{energy-based pointing game}.
Such checks have exposed shortcut learning~\cite{geirhos2020shortcut} in chest X-ray models: classifiers that rely on hospital-specific markers and laterality tokens~\cite{zech2018variable}
or on image regions outside the lungs~\cite{degrave2021ai}.
We follow this line in Stage~3: Grad-CAM on both of our models, and the energy share inside a lung mask instead of a bounding box.
The masks come from a PSPNet~\cite{zhao2017pspnet} provided by TorchXRayVision~\cite{cohen2022torchxrayvision}, trained on the ChestX-Det annotations~\cite{lian2021chestxdet}
of adult NIH chest X-rays~\cite{wang2017chestxray8}. Unlike those adult images, our dataset is paediatric, so we check the masks by eye.


% =====================================================================
% Dataset, preprocessing and baseline CNN  (author: Abhishek Singh, Stage 1)
% Every number below comes from results/ or data/splits/ in the repository:
%   results/stage1_eda/*.csv, data/splits/split_summary.json,
%   results/stage1_baseline/settings.json, results/stage1_metrics.csv
% Figures to upload to Overleaf (keep these paths):
%   figures/stage1/sample_grid.png          <- results/stage1_eda/sample_grid.png
%   figures/stage1/image_sizes.png          <- results/stage1_eda/image_sizes.png
%   figures/stage1/training_curves.png      <- results/stage1_baseline/training_curves.png
%   figures/stage1/confusion_matrices.png   <- results/stage1_baseline/confusion_matrices.png
% =====================================================================
\section{Dataset and Preprocessing}
\label{sec:data}

\subsection{Dataset}
\label{sec:data-dataset}
We use the \emph{Chest X-Ray Images (Pneumonia)} dataset~\cite{mooney2018kaggle}, the Kaggle release of the paediatric chest X-rays
collected by Kermany \etal~\cite{kermany2018identifying} at the Guangzhou Women and Children's Medical Center.
It contains 5,856 frontal chest X-rays in two classes: 1,583 \textsc{normal} and 4,273 \textsc{pneumonia}
(2,780 bacterial and 1,493 viral, which we treat as one class).
The images are split into official \texttt{train}, \texttt{val} and \texttt{test} folders (\cref{tab:data-splits}).
The download also contains a duplicate copy of all images in a nested \texttt{chest\_xray/} folder and a \texttt{\_\_MACOSX/} folder; we read only the top-level folders.
\Cref{fig:samples} shows example images.

\begin{figure*}[t]
  \centering
  \includegraphics[width=\linewidth]{figures/stage1/sample_grid.png}
  \caption{Random training images (fixed seed) with their original size. Text markers such as the letter \emph{R} appear in both classes.
           Pneumonia images are often wider relative to their height (\cref{sec:data-eda}).}
  \label{fig:samples}
\end{figure*}

\paragraph{Patients.}
Several images can belong to the same child (up to 30 for one patient), so images of one patient must never be spread across training and evaluation sets.
The filenames encode a patient identifier, from which we derive a patient ID:
\texttt{person\{N\}\_\{bacteria|virus\}} for pneumonia images and \texttt{IM-\{XXXX\}} or \texttt{NORMAL2-IM-\{XXXX\}} for normal images.
With this rule the official folders share no patients, and the dataset contains 4,097 patients (1,444 normal, 2,653 pneumonia).
If only \texttt{person\{N\}} were used, 170 identifiers would appear in both the official training and test folders.
Including the infection type removes all of this overlap, which strongly suggests that \texttt{person\{N\}\_bacteria} and \texttt{person\{N\}\_virus} are different children.
This is an inference from the file names and not a fact documented by the dataset authors (\cref{sec:data-limits}).

\subsection{Data splits}
\label{sec:data-splits}
The official validation folder contains only 16 images, which is too few for model selection.
We keep the official test folder unchanged as the test set and build a new validation set from the remaining 5,232 images
(official \texttt{train} plus \texttt{val}).
The split is \emph{patient-level} and stratified by class: for each class, the patients are shuffled with a fixed seed (42)
and assigned to validation until it holds 12.5\% of that class's images; all images of a patient go to the same set.
As a safety margin, the training/validation split groups pneumonia images by \texttt{person\{N\}} alone, keeping a person's bacterial and viral images on the same side,
so that the two sets share no identifier even if both types belonged to the same child.
The resulting sets (\cref{tab:data-splits}) share no patients (and training and validation share no \texttt{person\{N\}}), and both training and validation contain 74.2\% pneumonia images.
We first used scikit-learn's \texttt{StratifiedGroupKFold}, but with the same seed it produced different splits under different library versions.
We therefore implemented the split directly with NumPy's legacy \texttt{RandomState}, whose random stream is fixed, so the split is identical on every platform.
The split is stored as CSV files in the repository, and every stage of the project uses these exact files.

\begin{table}[t]
  \centering
  \small
  \begin{tabular}{lrrrr}
    \toprule
    Set & Images & Normal & Pneumonia & Patients \\
    \midrule
    \multicolumn{5}{l}{\emph{Official folders}} \\
    \texttt{train} & 5,216 & 1,341 & 3,875 & 3,654 \\
    \texttt{val}   &    16 &     8 &     8 &    15 \\
    \texttt{test}  &   624 &   234 &   390 &   428 \\
    \midrule
    \multicolumn{5}{l}{\emph{Our splits (used in all experiments)}} \\
    Train      & 4,578 & 1,180 & 3,398 (74.2\%) & 3,218 \\
    Validation &   654 &   169 &   485 (74.2\%) &   451 \\
    Test       &   624 &   234 &   390 (62.5\%) &   428 \\
    \bottomrule
  \end{tabular}
  \caption{Images and patients per set. Our validation set is a patient-level, stratified 12.5\% of the official \texttt{train} and \texttt{val} folders; the official test folder is unchanged.}
  \label{tab:data-splits}
\end{table}

\subsection{Exploratory analysis and potential shortcuts}
\label{sec:data-eda}
The images vary widely in size: widths range from 384 to 2,916 pixels and heights from 127 to 2,713 pixels.
Two properties of the raw files differ systematically between the classes (\cref{fig:sizes}):
\begin{itemize}
  \item \textbf{Size and aspect ratio.} Normal images are larger (median $1654\times1323$ pixels) than pneumonia images (median $1160\times776$),
        and their median width-to-height ratio is about 1.2, compared with about 1.5 for pneumonia images.
        A plain resize to a square input therefore squeezes the two classes by different amounts, so a model could learn the distortion instead of the lung findings.
  \item \textbf{Colour mode.} All 283 RGB files in the dataset are pneumonia images; every normal image is stored as single-channel grayscale.
\end{itemize}
Both properties are potential shortcuts that exist before any model is trained.
We remove the colour-mode cue by converting every image to grayscale (\cref{sec:data-preproc}).
We keep the plain resize: padding to a square would encode the aspect ratio in the width of the black bars and add artificial borders,
and cropping could cut off the edges of the lungs.
We therefore treat the aspect ratio as a known confound when interpreting the results.

\begin{figure}[t]
  \centering
  \includegraphics[width=\linewidth]{figures/stage1/image_sizes.png}
  \caption{Image size (left) and aspect ratio (right) per class. Pneumonia images are smaller and relatively wider, so a square resize distorts the two classes differently.}
  \label{fig:sizes}
\end{figure}

\subsection{Preprocessing and augmentation}
\label{sec:data-preproc}
All models in this project share the same preprocessing.
Each image is converted to grayscale, copied to three channels (so that ImageNet-pretrained networks can be used in later stages),
resized to $224\times224$ pixels without padding or cropping, and normalised with the ImageNet mean and standard deviation.
During training only, images are rotated by a random angle of up to $\pm10^\circ$.
We deliberately do not use horizontal flips: flipping a chest X-ray moves the heart to the other side and mirrors the laterality markers,
which are among the shortcuts this project studies.
The classes are imbalanced (74.2\% pneumonia in the training set).
We compensate in the loss with a positive-class weight equal to the ratio of normal to pneumonia images in the training set, $w_{+} = 1180/3398 = 0.347$.
Data loading, shuffling and augmentation are seeded, so training runs are reproducible.

\section{Baseline CNN}
\label{sec:baseline}

\paragraph{Architecture.}
As a baseline we train a custom CNN from scratch.
It has five convolutional blocks with 32, 64, 128, 256 and 256 channels.
Each block applies two $3\times3$ convolutions, each followed by batch normalisation~\cite{ioffe2015batch} and a ReLU, and then $2\times2$ max pooling.
The blocks are followed by global average pooling, dropout ($p=0.3$) and a single linear output: the logit of the pneumonia class.
The network has 2.35\,M parameters.
The last convolutional layer produces a $14\times14$ feature map, which Grad-CAM uses to produce the heatmaps in Stage~3.

\paragraph{Training.}
The model is trained with binary cross-entropy on the logit, using the positive-class weight from \cref{sec:data-preproc},
and the Adam optimiser~\cite{kingma2015adam} (learning rate $10^{-3}$, weight decay $10^{-4}$, batch size 32).
The learning rate is halved when the validation loss has not improved for two epochs.
Training stops when the validation loss has not improved for five epochs (at most 30 epochs), and we keep the weights with the lowest validation loss.
Model selection uses only the validation set: the test set is evaluated once, with the final model and a fixed decision threshold of 0.5.
The best validation loss was reached at epoch 3, and training stopped after epoch 8 (\cref{fig:curves}).
Training took about two minutes on a single NVIDIA T4 GPU with mixed precision, using PyTorch~\cite{paszke2019pytorch}.

\begin{figure}[t]
  \centering
  \includegraphics[width=\linewidth]{figures/stage1/training_curves.png}
  \caption{Baseline training: loss (left) and validation AUC and F1 (right) per epoch. The weights from epoch 3 (lowest validation loss) are kept.}
  \label{fig:curves}
\end{figure}

\paragraph{Results.}
\Cref{tab:baseline} reports the baseline on the validation and test sets, with pneumonia as the positive class.
On validation, the model reaches 96.8\% accuracy and an AUC of 0.992.
On the test set, accuracy drops to 80.1\% (AUC 0.902).
The confusion matrices (\cref{fig:cm}) show where the difference comes from.
The model misses only 7 of the 390 pneumonia test images (recall 0.982), but it labels 117 of the 234 normal test images as pneumonia,
so it correctly identifies only 50.0\% of normal test images, compared with 95.9\% on validation.
Our validation set is drawn from the same pool as the training data, while the official test set appears to differ in distribution,
and the model's errors on it are almost entirely false alarms on normal images.
We did not tune the threshold or any other setting on the test set.
DenseNet-121 later raises test accuracy to 0.873 and AUC to 0.971, and reduces false positives on normal images from 117 to 78 (\cref{sec:densenet-results}).
This gap is a starting point for the attention analysis: if the false positives are driven by regions outside the lungs,
the Grad-CAM metric and the masking experiment should reveal it.

\begin{table}[t]
  \centering
  \small
  \begin{tabular}{lccccc}
    \toprule
    Set & Acc. & Prec. & Recall & F1 & AUC \\
    \midrule
    Validation & 0.968 & 0.985 & 0.971 & 0.978 & 0.992 \\
    Test       & 0.801 & 0.766 & 0.982 & 0.861 & 0.902 \\
    \bottomrule
  \end{tabular}
  \caption{Custom CNN baseline (threshold 0.5, pneumonia = positive class). Source: \texttt{results/stage1\_metrics.csv}.}
  \label{tab:baseline}
\end{table}

\begin{figure}[t]
  \centering
  \includegraphics[width=\linewidth]{figures/stage1/confusion_matrices.png}
  \caption{Confusion matrices of the baseline on validation (left) and test (right). On the test set, 117 of 234 normal images are predicted as pneumonia.}
  \label{fig:cm}
\end{figure}

\subsection{Limitations of the data}
\label{sec:data-limits}
% Can be moved into the shared Limitations / Discussion section.
The patient IDs are inferred from file names (\cref{sec:data-dataset}); zero overlap between the official folders is strong evidence, but not proof, that our split separates patients.
The aspect-ratio difference between the classes (\cref{sec:data-eda}) is not removed by our preprocessing, and lung masking (Stage~4) does not remove it either, because the masked lungs are still distorted.
Finally, all images come from a single medical centre, so none of our test results measure how well the models generalise to other hospitals.

% =====================================================================
% DenseNet-121 method, training setup and comparison  (author: Ajit Yadav, Stage 2)
% Every number below comes from results/ in the repository:
%   results/stage2_metrics.csv, results/stage2_comparison.csv, results/stage1_metrics.csv,
%   results/stage2_phaseA/settings.json + val_metrics.csv, results/stage2_phaseB/tuning.csv,
%   results/stage2_densenet/settings.json
% Figures to upload to Overleaf (keep these paths):
%   figures/stage2/confusion_matrices.png   <- results/stage2_densenet/confusion_matrices.png
%   figures/stage2/roc.png                  <- results/stage2_densenet/roc.png
% Labels used from Abhishek's section: sec:data-splits, sec:data-preproc, sec:baseline, tab:baseline, sec:data-limits
% =====================================================================
\section{DenseNet-121 Transfer Learning}
\label{sec:densenet}

\subsection{Method}
\label{sec:densenet-method}
The stronger model is DenseNet-121~\cite{huang2017densenet} initialised with ImageNet~\cite{deng2009imagenet} weights, following CheXNet~\cite{rajpurkar2017chexnet}.
We replace the original 1000-class layer by a dropout layer (rate 0.2) and one linear layer with a single output logit, where the sigmoid of the logit is the probability of pneumonia.
The network has 6.95\,M parameters.
The data pipeline is exactly the one used for the baseline (\cref{sec:data-preproc}): grayscale images copied to three channels, plain resize to $224\times224$,
ImageNet mean and standard deviation, and a train-only rotation of $\pm10^\circ$ without horizontal flip.
It also uses the same patient-level splits (\cref{sec:data-splits}), the same weighted loss, and the same metrics as the baseline, so the two models differ only in the architecture and its initialisation.

\subsection{Training setup}
\label{sec:densenet-training}
We train in two phases, as in CheXNet-style transfer learning.
In \emph{Phase~A}, the backbone is frozen (its BatchNorm layers are kept in evaluation mode, so their statistics do not change) and only the new output layer, 1,025 parameters, is trained
with Adam (learning rate $10^{-3}$, weight decay $10^{-4}$, batch size 32) for up to 20 epochs.
The best epoch (19) has a validation loss of 0.065 and a validation AUC of 0.990.
In \emph{Phase~B}, we start from the Phase~A weights, unfreeze the whole network and fine-tune all 6.95\,M parameters with a low learning rate
(Adam, weight decay $10^{-4}$, for at most 15 epochs).
Both phases minimise \texttt{BCEWithLogitsLoss} with the class weight computed on the training split only ($\text{pos\_weight}=0.347$),
halve the learning rate when the validation loss stops improving (patience 2), and stop early on the validation loss (patience 5 in Phase~A, 4 in Phase~B).
Mixed precision is used on the GPU, and all runs use seed 42.

\paragraph{Tuning.}
We tuned the learning rate and the batch size of Phase~B with three runs (\cref{tab:tuning}).
The model is selected by the lowest \emph{validation} loss; the test set is not used for any choice and is evaluated once, for the selected model only.
The selected run uses a learning rate of $10^{-4}$ and batch size 16 (best epoch 3).

\begin{table}[t]
  \centering
  \small
  \begin{tabular}{lrrrr}
    \toprule
    Run (Phase B) & LR & Batch & Val loss & Val AUC \\
    \midrule
    \textbf{lr1e-4\_bs16} & $10^{-4}$ & 16 & \textbf{0.0139} & 0.9995 \\
    lr1e-4                & $10^{-4}$ & 32 & 0.0206 & 0.9994 \\
    lr5e-5                & $5\cdot10^{-5}$ & 32 & 0.0278 & 0.9984 \\
    \bottomrule
  \end{tabular}
  \caption{Phase~B tuning on the validation set (patient-level split). The run with the lowest validation loss (bold) is the final model.}
  \label{tab:tuning}
\end{table}

\subsection{Results and comparison with the baseline}
\label{sec:densenet-results}
\Cref{tab:comparison} compares the final DenseNet-121 with the baseline CNN (\cref{sec:baseline}) on the same official test set with the same metrics (threshold 0.5).
DenseNet-121 raises accuracy from 0.801 to 0.873, F1 from 0.861 to 0.908 and AUC from 0.902 to 0.971.
The number of normal images predicted as pneumonia falls from 117 to 78 of 234, and only one of 390 pneumonia images is missed (baseline: 7).
The pretrained backbone provides better features than a network trained from scratch on about 4,500 images, which is the usual benefit of transfer learning.
\Cref{fig:densenet-cm} shows the confusion matrices and \cref{fig:densenet-roc} the ROC curves.

\begin{table}[t]
  \centering
  \small
  \begin{tabular}{lrrrrrrr}
    \toprule
    Model & Acc. & Prec. & Rec. & F1 & AUC & FP & FN \\
    \midrule
    Custom CNN   & 0.801 & 0.766 & 0.982 & 0.861 & 0.902 & 117 & 7 \\
    DenseNet-121 & \textbf{0.873} & \textbf{0.833} & \textbf{0.997} & \textbf{0.908} & \textbf{0.971} & \textbf{78} & \textbf{1} \\
    \bottomrule
  \end{tabular}
  \caption{Test-set results of the baseline CNN and DenseNet-121 (unmasked images, threshold 0.5; 234 normal and 390 pneumonia images). FP: normal images predicted as pneumonia; FN: missed pneumonia images.}
  \label{tab:comparison}
\end{table}

\begin{figure}[t]
  \centering
  \includegraphics[width=\linewidth]{figures/stage2/confusion_matrices.png}
  \caption{Confusion matrices of DenseNet-121 on the validation and test sets (threshold 0.5).}
  \label{fig:densenet-cm}
\end{figure}

\begin{figure}[t]
  \centering
  \includegraphics[width=0.7\linewidth]{figures/stage2/roc.png}
  \caption{ROC curves of DenseNet-121 on the validation and test sets.}
  \label{fig:densenet-roc}
\end{figure}

\paragraph{Validation--test gap.}
Like the baseline, DenseNet-121 performs much better on the validation set (accuracy 0.985, AUC 0.9995) than on the test set (accuracy 0.873, AUC 0.971).
The validation set is drawn from the same distribution as the training data, whereas the official test set differs from it, so the test metrics are the ones that matter.
DenseNet-121 narrows the gap (0.111 accuracy difference, baseline: 0.167) but does not remove it, and 78 normal test images are still predicted as pneumonia.
Whether these errors come from image regions outside the lungs is examined with Grad-CAM in the next section. % TODO: link to Aman's Grad-CAM section (Stage 3)

\paragraph{Limitations.}
All images come from one medical centre (\cref{sec:data-limits}), the three tuning runs cover only a small part of the hyperparameter space,
and all numbers come from a single training run with seed 42, so small differences between models should not be over-interpreted.
% TODO (Himanshu, Stage 4): add the lung-masked DenseNet-121 row to \cref{tab:comparison} and the attention results.

% =====================================================================
% Grad-CAM, lung masks, attention metric and shortcut findings  (author: Aman Singh, Stage 3)
% Every number below comes from results/ in the repository:
%   results/stage3_metrics.csv, results/stage3/attention_summary.csv, results/stage3/attention_robustness.csv,
%   results/stage3/lung_mask_summary.csv, results/stage3/sample_images.csv, results/stage3/settings.json
% Figures to upload to Overleaf (keep these paths):
%   figures/stage3/lung_masks_random.png          <- results/stage3/lung_masks_random.png
%   figures/stage3/attention_boxplot.png          <- results/stage3/attention_boxplot.png
%   figures/stage3/false_positives_densenet121.png <- results/stage3/false_positives_densenet121.png
% Labels used from other sections: sec:data-preproc, sec:data-eda, sec:densenet, tab:comparison
% Citation keys: selvaraju2017gradcam, geirhos2020shortcut, zech2018variable, degrave2021ai (references_abhishek.bib);
%   zhao2017pspnet, lian2021chestxdet, wang2017chestxray8, cohen2022torchxrayvision, wang2020scorecam (references_aman.bib)
% =====================================================================
\section{Where Do the Models Look? Grad-CAM and Lung Attention}
\label{sec:gradcam}

Both classifiers reach a high test AUC (\cref{tab:comparison}), but a high score does not show that a model uses the lungs.
In this section we measure where the custom CNN and DenseNet-121 look, and how much of that attention falls on the lungs.

\subsection{Grad-CAM for a single-logit classifier}
\label{sec:gradcam-method}
Grad-CAM~\cite{selvaraju2017gradcam} explains a decision with the feature maps $A^k$ ($k=1,\dots,K$) of a late convolutional layer.
Each map is weighted by the average gradient of the explained score $s$, and the weighted sum is passed through a ReLU:
\begin{equation}
  \alpha_k = \frac{1}{hw}\sum_{i,j}\frac{\partial s}{\partial A^k_{ij}}, \qquad
  L = \mathrm{ReLU}\Big(\sum_k \alpha_k A^k\Big).
  \label{eq:gradcam}
\end{equation}
Our models have one output logit $z$ with $\sigma(z)=P(\text{pneumonia})$ and no separate NORMAL output.
We therefore explain the \emph{predicted} class: $s=z$ when the image is predicted PNEUMONIA ($\sigma(z)\ge0.5$) and $s=-z$ when it is predicted NORMAL,
so every heatmap shows the evidence for the class the model chose.
We hook the activations after the last ReLU, which are the features the classifier actually receives:
\texttt{features.block5.relu2} for the custom CNN ($256\times14\times14$) and the post-ReLU feature map of DenseNet-121 ($1024\times7\times7$).
The map $L$ is upsampled bilinearly to $224\times224$, the input size, so it lines up pixel by pixel with the lung masks.
A toy model with a known answer confirms that the heatmap lands on the right image region for both predicted classes.

\subsection{Lung masks}
\label{sec:gradcam-masks}
We segment the lungs of all 5,856 images with the pretrained PSPNet~\cite{zhao2017pspnet} from TorchXRayVision~\cite{cohen2022torchxrayvision},
which was trained on the ChestX-Det annotations~\cite{lian2021chestxdet} of adult NIH chest X-rays~\cite{wang2017chestxray8}.
We join its left- and right-lung outputs into one mask.
The classifiers see the whole image squashed to $224\times224$ (\cref{sec:data-preproc}), so we also resize the whole image without cropping before segmentation,
and resize the mask back to $224\times224$ with nearest-neighbour interpolation to keep it binary.
\Cref{fig:masks} shows masks for 24 random images.
% TODO (Aman): after the eye check of results/stage3/lung_masks_random.png, state how many of the 24 masks look right
% (the failures seen so far are pneumonia images with a partly missed, consolidated lung).

The segmentation model was not trained on children, and it fails on some pneumonia images: a lung that is white from consolidation is partly or entirely missed.
We flag a mask when the lungs cover less than 8\% or more than 60\% of the image, or when one lung is under 5\% of the image or under 40\% of the other lung.
Of the 5,856 masks, 226 are flagged; 225 of them are PNEUMONIA images (181 train, 25 validation, 19 test).

\begin{figure*}[t]
  \centering
  \includegraphics[width=\linewidth]{figures/stage3/lung_masks_random.png}
  \caption{Lung masks (red, outlined in green) for 24 random images, drawn on the $224\times224$ input the classifiers see. Most masks cover both lung fields; failures occur on pneumonia images where consolidated lung is missed.}
  \label{fig:masks}
\end{figure*}

\subsection{Lung-attention metric}
\label{sec:gradcam-metric}
For a heatmap $L\ge0$ and a lung mask $M\in\{0,1\}$ of the same size, the lung attention is the share of the heatmap inside the lungs:
\begin{equation}
  \text{attention\_in\_lung} = \frac{\sum_{i,j} L_{ij} M_{ij}}{\sum_{i,j} L_{ij}}.
  \label{eq:attention}
\end{equation}
This is the energy-based pointing game of Score-CAM~\cite{wang2020scorecam} with the lung mask in place of a bounding box.
The value alone is hard to read: a heatmap spread evenly over the image would score the lung area $\frac{1}{hw}\sum_{i,j}M_{ij}$, about 38\% on our test images.
We therefore report the lung area as the chance level next to each result, and the \emph{lift}, the ratio of the two, where a lift above 1 means more attention on the lungs than chance.
We also report the \emph{border share}, the part of the heatmap in the outer 10\% of the image on each side, where text markers and scanner labels sit;
an even heatmap has a border share of $1-0.8^2=36\%$.
All values are computed on the 624 test images for both models.

\subsection{Results}
\label{sec:gradcam-results}
\begin{table}[t]
  \centering
  \caption{Lung attention on the test set (mean over images). Lung area is the chance level: the score of an evenly spread heatmap. Lift $>1$ is the share of images with more attention on the lungs than chance.}
  \label{tab:attention}
  \small
  \begin{tabular}{llrrrr}
    \toprule
    Model & Images & Attention & Lung area & Lift $>1$ & Border \\
    \midrule
    Custom CNN   & all (624)       & 43.4\% & 37.7\% & 76.8\% & 26.8\% \\
                 & NORMAL (234)    & 42.3\% & 39.4\% & 64.5\% & 27.7\% \\
                 & PNEUMONIA (390) & 44.1\% & 36.7\% & 84.1\% & 26.2\% \\
    \midrule
    DenseNet-121 & all (624)       & 39.5\% & 37.7\% & 60.1\% & 30.9\% \\
                 & NORMAL (234)    & 42.9\% & 39.4\% & 68.8\% & 28.5\% \\
                 & PNEUMONIA (390) & 37.5\% & 36.7\% & 54.9\% & 32.4\% \\
    \bottomrule
  \end{tabular}
\end{table}

\Cref{tab:attention} shows that both models put only slightly more of their heatmap on the lungs than an evenly spread heatmap would:
43.4\% for the custom CNN and 39.5\% for DenseNet-121, against a lung area of 37.7\%.
The more accurate model is the one closer to chance.
On PNEUMONIA images, DenseNet-121 puts 37.5\% of its attention on the lungs, which cover 36.7\% of the image, and only 54.9\% of these images have a lift above~1.
High accuracy therefore does not mean that the model bases its decision on the lungs.
\Cref{fig:attention-box} shows the spread per class and for correct and wrong predictions.

\begin{figure*}[t]
  \centering
  \includegraphics[width=\linewidth]{figures/stage3/attention_boxplot.png}
  \caption{Lung attention per test image, by true class and correct or wrong prediction. The red dash is the mean lung area of each group, the chance level.}
  \label{fig:attention-box}
\end{figure*}

\subsection{Shortcut examples}
\label{sec:gradcam-shortcuts}
The clearest shortcut appears in DenseNet-121's most confident mistakes (\cref{fig:shortcuts}).
For NORMAL test images that it predicts as PNEUMONIA with probability 1.00, the heatmap concentrates on text burned into the image:
the \texttt{KV58}/\texttt{KV59}/\texttt{KV60} exposure labels in the bottom-right corner and the timestamps in the top-right corner, not the lungs.
The numbers agree with the pictures. On NORMAL images that are wrongly predicted as PNEUMONIA, the border share is 35.8\% for the custom CNN (19.6\% on correctly predicted NORMAL images)
and 34.6\% for DenseNet-121 (25.4\% on correct ones).
The custom CNN's heatmaps often sit on the diaphragm and the strip below the lungs.
These findings match the shortcuts reported for chest X-ray classifiers by Zech~\etal~\cite{zech2018variable} and DeGrave~\etal~\cite{degrave2021ai},
and the general pattern of shortcut learning described by Geirhos~\etal~\cite{geirhos2020shortcut}. They also add to the recording differences between the classes found in the EDA (\cref{sec:data-eda}).
% TODO (Aman): add the 6-8 hand-picked examples of Stage 3.4 (issue #14) from results/stage3/ and describe each one.

\begin{figure*}[t]
  \centering
  \includegraphics[width=\linewidth]{figures/stage3/false_positives_densenet121.png}
  \caption{DenseNet-121's most confident false positives (NORMAL test images predicted as PNEUMONIA). The heatmaps focus on burned-in text such as the \texttt{KV58}/\texttt{KV59}/\texttt{KV60} exposure labels (bottom right) and timestamps (top right) instead of the lungs. Green: lung mask outline.}
  \label{fig:shortcuts}
\end{figure*}

\subsection{Limitations}
\label{sec:gradcam-limits}
The lung masks come from a model trained on adult X-rays and can miss consolidated lung, so attention on diseased lung may be counted as outside the lungs.
This affects PNEUMONIA images only. Leaving out the 19 flagged test images changes the mean lung attention by less than 0.01 for every model and class (largest change: custom CNN on PNEUMONIA images, 44.1\% to 44.7\%),
but masks that are only partly wrong are not all flagged.
DenseNet-121's heatmaps are computed on a $7\times7$ map (the custom CNN's on $14\times14$), and coarser maps spread out more after upsampling, which moves the lung attention towards the chance level.
Grad-CAM shows which regions raised the score, not why, and it does not prove that the model would fail without them;
\cref{sec:masking} tests this directly by training on lung-masked images. % TODO (Himanshu): check the label of the Stage 4 section

% Stage 4 (author: Himanshu Kumar).
% Training numbers come from results/stage4_masked_densenet/summary.json and
% results/stage2_phaseA/settings.json / results/stage2_phaseB/lr1e-4_bs16/settings.json.
% Test numbers come from results/stage4_metrics.csv,
% results/stage4_cross_test_metrics.csv, results/stage4_comparison.csv,
% results/stage4_masked_densenet/masked_test_predictions_attention.csv,
% results/stage2_densenet/test_predictions.csv and results/stage3/lung_mask_summary.csv.
\section{Lung-Masked Training and Cross-Test}
\label{sec:masking}

\subsection{Training protocol}
We train a second DenseNet-121 after masking pixels outside the lungs. Each X-ray is converted to grayscale,
resized to $224\times224$, and set to black outside the binary lung mask. The same committed patient-level
splits and Stage~2 training settings are used. Phase~A trains the classifier head and Phase~B fine-tunes the
whole network; the checkpoint is selected by validation loss, without evaluating the test set.

The selected masked model reached its lowest Phase~A validation loss of 0.083285 at epoch~19
(validation AUC 0.986409), then its lowest Phase~B validation loss of 0.028584 at epoch~10
(validation AUC 0.998512). As in Stage~2, Phase~A used batch size 32 and learning rate $10^{-3}$;
Phase~B used batch size 16 and learning rate $10^{-4}$. Both phases used seed 42 and the same
threshold and preprocessing settings as Stage~2.

\subsection{Test evaluation}
After selection, we evaluate the masked model on masked test images and cross-test the Stage~2 unmasked
checkpoint on those same masked images. For the masked model, the same forward/backward pass supplies both
predictions and Grad-CAM maps. We report the share of heatmap energy inside the lung masks with
\texttt{attention\_in\_lung}, shown as a percentage, alongside accuracy, recall, F1 and AUC.

\begin{table}[t]
  \centering
  \small
  \begin{tabular}{llrrrrr}
    \toprule
    Model & Training data & Acc. & Rec. & F1 & AUC & Lung attention (\%) \\
    \midrule
    Custom CNN   & Unmasked & 0.8013 & 0.9821 & 0.8607 & 0.9016 & 43.41\% \\
    DenseNet-121 & Unmasked & 0.8734 & 0.9974 & 0.9078 & 0.9713 & 39.52\% \\
    DenseNet-121 & Lung-masked & 0.8702 & 0.9949 & 0.9055 & 0.9759 & 51.79\% \\
    \bottomrule
  \end{tabular}
  \caption{Final test comparison on the committed 624-image test split. Baseline and unmasked DenseNet-121 values are from Stage~1 and Stage~3; the masked-model values are from the Stage~4.2/4.3 Kaggle evaluation.}
  \label{tab:masking-comparison}
\end{table}

\paragraph{Cross-test.}
The Stage~2 unmasked checkpoint evaluated on masked test images achieves 0.6314 accuracy, 0.6290 precision,
1.0000 recall, 0.7723 F1 and 0.8625 AUC. It predicts 230 of the 234 NORMAL images as PNEUMONIA
(TN=4, FP=230, FN=0, TP=390). This is a substantial drop from its matched unmasked test result
(0.8734 accuracy and 0.9713 AUC), showing sensitivity to removal of extra-pulmonary image context.
The masked inputs also introduce black backgrounds that the unmasked model did not see during training,
so this cross-test drop may partly reflect an input-distribution change.

\paragraph{Interpretation and limitations.}
On matched inputs, the masked model is close to the unmasked DenseNet-121 in accuracy (0.8702 vs.~0.8734),
recall (0.9949 vs.~0.9974) and F1 (0.9055 vs.~0.9078), with a slightly higher AUC (0.9759 vs.~0.9713).
Its mean lung attention is 51.79\%, compared with 39.52\% for the unmasked model and a mean lung-area
baseline of 37.70\%; the masked model's mean attention-to-area lift is 1.40. The masked model has 79 false
positives versus 78 for the unmasked model, with 68 images in common. Both highlighted Stage~3 examples
(NORMAL2-IM-0252 and NORMAL2-IM-0221) remain false positives after masking. Their corner text shows where
the unmasked model attends, but does not explain most errors; the test-set gap remains. The cross-test
indicates sensitivity to changing pixels outside the lungs, but does not identify a specific shortcut or
establish causation. The Stage~3 segmentation caveat also applies: 19 test masks and 182 training masks are
flagged for possible segmentation failure, and all remain included. The eye check also found partial mask
misses that were not flagged. Grad-CAM attention should therefore be interpreted alongside mask quality
rather than as a definitive measure of causal model reliance.


\section{Conclusion}
The baseline and DenseNet-121 achieve strong test performance, but their Grad-CAM maps and selected examples show that high classification scores alone do not establish reliance on pulmonary evidence. On matched test inputs, masked DenseNet-121 retains similar classification performance and assigns more Grad-CAM attention to the lung masks. Its 79 false positives overlap with 68 of the unmasked model's 78 false positives, and both highlighted corner-text examples remain misclassified after masking; those examples show attention location but do not explain most errors. The unmasked checkpoint's cross-test performance falls substantially on masked inputs, which introduce black backgrounds it did not see during training. This indicates sensitivity to the changed input but does not isolate the effect of removing extra-pulmonary features. The adult-trained segmentation model's errors, including 19 flagged test masks and 182 flagged training masks, as well as partial misses not caught by the flag and the single-centre dataset, constrain these conclusions. Future work should validate on external cohorts and with more reliable lung annotations.

\bibliographystyle{plain}
\bibliography{references_abhishek,references_ajit,references_aman}

\end{document}
